\begin{textsample}{NEG Dim 7 – human – Score -2.00 – mg/mg\_ef\_linguagens\_ed\_fis\_13\_p4.txt}  \label{ex:f7_neg_019}
Chamada a atenção para os limites a serem superados, ressaltamos a importância da construção de práticas avaliativas na Educação Física escolar baseadas em múltiplas estratégias, procedimentos, \textbf{instrumentos} e recursos que possam nortear a adoção e revisão de metodologias e ações pedagógicas que garantam o desenvolvimento das habilidades e competências propostas para a etapa do ensino fundamental.
Nesse sentido, mensurações quantitativas e qualitativas têm espaço e valor, desde que contextualizadas e vinculadas aos objetivos almejados, sempre em prol dos princípios aqui propostos : cooperação, inclusão, solidariedade, respeito e empatia.

% matched lemmas: instrumento
\end{textsample}
